\section{Limitations}
\label{sec:limitations}

\subsection{What we do not claim}

\begin{center}
\begin{tabular}{p{0.42\textwidth}p{0.52\textwidth}}
\toprule
Claim we avoid & Why \\
\midrule
``feedback is useless'' / ``we established a zero effect''
  & No equivalence margin was preregistered. We use only the wording of
    Section~\ref{sec:c3}. \\
``exactly half of \gate{C3} was an acceptance artifact''
  & The reduction differs by regime (62\%, 42\%) and $n = 3$. We use only the wording
    of Section~\ref{sec:ablation}. \\
``reinforcement learning fails at optimizer control''
  & We did not train a policy. \\
``the deterministic setting is $N$ times better than the stochastic one''
  & GE does not match FLOPs across batch sizes (Section~\ref{sec:ge}). \\
``headroom decreases with condition number''
  & We observed non-monotonicity only. We claim monotonicity in neither direction. \\
``the method fails on nonlinear problems''
  & Both Rosenbrock variants were ruled unusable as benchmark defects. \\
``the micro-neural results overturn the quadratic results''
  & $n = 3$ exploratory versus $n = 40$ held-out. \\
Any value formed by adding or subtracting paired medians
  & The median is not linear. Comparisons use direct paired measurements only
    (Section~\ref{sec:no-subtraction}). \\
\bottomrule
\end{tabular}
\end{center}

\subsection{Protocol deviations}
\label{sec:deviations}

\begin{description}
  \item[E1] The configuration selection rule was finalized after inspecting the gate
    table of a single-instance dry run (Section~\ref{sec:deviation-selection}).
  \item[E2] (quad $\kappa{=}10^5$, seed 2) overlaps between the initial dev audit and
    configuration selection (Section~\ref{sec:seeds}).
  \item[E3] Some executions ran with uncommitted changes, recorded as
    \texttt{code\_dirty}.
  \item[E4] Tests ran concurrently with the held-out execution. Wall-clock records are
    therefore inaccurate, but numerical results and verdicts are unaffected
    (single-threaded deterministic arithmetic).
  \item[E5] Summary-level provenance omitted the commit field because of a reporting
    bug; the corresponding per-run raw records retained the commit, allowing
    deterministic reconstruction.
  \item[E6] The acceptance-criterion ablation is not a single-factor change
    (Section~\ref{sec:ablation-limits}).
  \item[E7] The R1 column of the alternative criterion is a budget perturbation
    (Section~\ref{sec:ablation-limits}).
  \item[E8] Gate \gate{C1} uses a single-seed diagnostic baseline and was not used for
    any verdict.
  \item[E9] Three-layer reporting is not applied to gates \gate{A1} and \gate{B},
    which use a single statistic.
  \item[E10] An earlier diagnostic script used an upper-middle convention for even
    sample counts, whereas the main paired analysis used the conventional arithmetic
    median. All reporting code was unified before manuscript generation; the selected
    configuration and qualitative conclusions were unchanged.
  \item[E11] Run identifiers initially embedded the full target table rather than only
    the targets used by the corresponding task family. Adding a target entry for an
    unrelated task family therefore invalidated the cost-to-target runs of the
    quadratic suite. We restricted the identifier to the targets actually used and
    re-executed those runs; the fixed-budget comparisons were unaffected, and the
    cost-to-target statistic is a re-aggregation of the same measurements.
  \item[E12] The protocol specified that a tag be created at the
    configuration-freeze point. The tag exists in the private repository, but it was
    created after the confirmatory runs had completed and it points at the final
    Stage~2 commit rather than at the freeze commit. It therefore cannot attest that
    the configuration was fixed before the held-out runs; that ordering is attested by
    the dated decision record and by the freeze commit itself. We did not move the tag
    retroactively. The tag is not part of the public export: the commit it marks
    precedes the removal of device names from committed artifacts, and the public
    repository does not share history with the private one. It is, however, present on
    the private repository's remote, so withholding it from the export limits the
    exposure to that repository rather than eliminating it.
\end{description}

\subsection{Scope and confounds}
\label{sec:scope}

\begin{description}
  \item[L1] The primary confirmatory results are limited to synthetic ill-conditioned
    SPD quadratics: $d = 100$ fixed, four values of $\kappa$, ten held-out seeds.
  \item[L2] The micro-neural results are exploratory evidence at $n = 3$ per regime.
    We do not quote CIs or $p$-values.
  \item[L3] The acceptance ablation includes the evaluation forward cost and is
    therefore not a single-factor change; the acceptance criterion and the number of
    effective steps change together.
  \item[L4] GE matches oracle calls within a regime, not FLOPs across batch sizes.
    Cross-regime absolute comparisons are descriptive.
  \item[L5] The fixed-evaluation criterion is not a relaxation. Measured rejection
    rates rose under it; it is the stricter rule.
  \item[L6] Beam-search cost is far larger than the object-level budget:
    $194{,}095 \GE / 150 \GE = 1{,}294\times$.
  \item[L7] The selection rule was fixed before the full execution but after
    inspecting the gate table of a cost-measurement dry run. Recorded as a protocol
    deviation (E1).
  \item[L8] We did not run PPO, so we can draw no conclusion about learned-policy
    performance.
  \item[L9] The results do not generalize to Hessian-free Newton methods as a family,
    nor to neural optimization in general.
\end{description}

The Rosenbrock benchmark defect and the median-convention correction appear in
Section~\ref{sec:reproducibility} and in the deviation list above and are not
repeated here.
